\documentclass[10pt, letterpaper]{article}

\usepackage{../../../../template/template}

% ------------------------
% Impostazioni documento
% ------------------------

\setDocTitle{Verbale esterno 08/04/2026}
\setDocVersion{-}
\setDocData{08/04/2026}
\setDocIndex{ve\_2026\_04\_08}
\setDocState{2}
\setDocRecipient{2}

\begin{document}

\makeTitlePage

\newpage

\tableofcontents

\newpage

% -------------------------------------------------
% Informazioni riunione
% -------------------------------------------------
\makeMeetingInfoTable{08/04/2026}{16:00}{17:30}{in presenza}

% ------------------------
% Partecipanti
% ------------------------
\addParticipant{Filippo}{Venzo}{Programmatore}{P}
\addParticipant{Valeria}{Baleanu}{Programmatore}{P}
\addParticipant{Luca}{Granziero}{Amministratore}{P}
\addParticipant{Christian}{Libralato}{Responsabile}{P}
\addParticipant{Francesco}{Pasqual}{Programmatore}{P}
\addParticipant{Leonardo}{Pellizzon}{Programmatore}{P}
\addParticipant{Giuseppe}{De Fina}{Verificatore}{P}

\makeMeetingParticipantsTable

Oltre a SnakeByte hanno partecipato: responsabili aziendali, team Cloud, team UX/UI.
% -------------------------------------------------
% Ordine del giorno
% -------------------------------------------------
\section{Ordine del giorno}
\begin{itemize}
      \item Presentazione finale del progetto View4Life;
      \item illustrazione delle funzionalità principali dell'applicativo;
      \item approfondimento sull'architettura tecnica e scelte tecnologiche;
      \item demo applicativa e simulazione flussi operativi;
      \item approvazione MVP.
      \item feedback conclusivi e retrospettiva.
\end{itemize}

% -------------------------------------------------
% Approfondimento
% -------------------------------------------------
\section{Approfondimento}

\subsection*{Presentazione finale del progetto View4Life}
Il gruppo ha introdotto il progetto View4Life, delineandone gli obiettivi principali: lo sviluppo di un'applicazione web per il monitoraggio di residenze protette, la gestione di dispositivi IoT e l'analisi dei consumi energetici. È stata ripercorsa la cronologia del progetto, strutturata in due sprint di progettazione seguiti dallo sviluppo tecnico e dalla redazione documentale.

\subsection*{Illustrazione delle funzionalità principali dell'applicativo}
Sono state illustrate le caratteristiche core del sistema:
\begin{itemize}
    \item \textbf{Autenticazione}: Gestione multi-livello (Amministratore/OSS) con integrazione account \textit{myVimar} e gestione dei token.
    \item \textbf{Dashboard}: Panoramica centralizzata su allarmi, consumi e struttura topologica.
    \item \textbf{Gestione Allarmi}: Sistema di configurazione basato su regole personalizzabili (soglie, orari, priorità) e gestione in tempo reale tramite WebSocket.
    \item \textbf{Analytics}: Visualizzazione dati temporali e generazione di suggerimenti energetici tramite il modello LLM \textit{Groq$_G$}.
\end{itemize}

\subsection*{Approfondimento sull'architettura tecnica e scelte tecnologiche}
Il team ha approfondito lo stack tecnologico composto da Angular, NestJS, PostgreSQL e TimescaleDB. È stato spiegato l'uso di Docker per la containerizzazione e il deploy su droplet Digital Ocean. Particolare attenzione è stata posta sulla scelta di PostgreSQL per il caching e sull'integrazione con i servizi \textit{Vimar Cloud$_G$} per la sottoscrizione ai data point.

\subsection*{Demo applicativa e simulazione flussi operativi}
Durante la sessione pratica è stata mostrata l'interoperabilità del sistema:
\begin{itemize}
    \item Creazione e risoluzione di un allarme con notifica istantanea.
    \item Navigazione tra i reparti e assegnazione dinamica degli operatori.
    \item Invio di comandi ai dispositivi con feedback immediato in interfaccia.
    \item Visualizzazione dei grafici di analytics e dei suggerimenti generati dall'AI.
\end{itemize}

\subsection*{Approvazione MVP}
Al termine della demo la proponente ha approvato pienamente gli obiettivi raggiunti a livello di requisiti e di implementazioni, ritenendo anche l'organizzazione degna di nota. 
In data 09/04/2026 è anche prevista la presentazione dell'MVP al team Marketing Vimar, con particolare attenzione alla demo dell'applicativo.

\subsection*{Feedback conclusivi e retrospettiva}
I membri del gruppo hanno condiviso le lezioni apprese, evidenziando la crescita nelle metodologie Agile e la gestione delle criticità comunicative. 
In particolare, da una Retrospective Review finale è emerso un alto grado di soddisfazione dell'andamento del gruppo durante il progetto, sia a livello di obiettivi raggiunti che di comunicazione interna e 
gestione dei problemi, oltre che nel rispetto delle suddivisione dei ruoli e dei compiti. 
Ciò che invece ha generato più problemi è stato il rispetto delle scadenze, dato l'alto carico di lavoro dell'ultimo periodo, e la capacità di effettuare previsioni accurate sia a livello di scadenze, che di suddivisione dei ruoli.
L'esperienza è stata infine definita, dall'intero gruppo, molto formativa e anche utile per le prospettive future.


% -------------------------------------------------
% Decisioni
% -------------------------------------------------

% -------------------------------------------------
% Azioni
% -------------------------------------------------
\addTodo{\noOne}{Revisione finale dei deliverables, della documentazione tecnica e del manuale utente in base ai commenti ricevuti}{\teamName}{13/04/2026}
\addTodo{\noOne}{Consegnare i deliverables via email realizzando un documento di consultazione tramite puntatori.}{\teamName}{13/04/2026}

\makeTodoTable

\end{document}